\chapter{Objeto de estudio de la carrera}

\section{Definición del objeto de estudio}

La Matemática se entiende como conceptos y entes abstractos, sus propiedades y las relaciones entre estos. Estas entidades abstractas mayoritariamente responden  a versiones idealizadas de la realidad, aunque el proceso de transformación del objeto real observado al ente abstracto que lo representa puede ser intrincado, incluso podría llegar a tomar siglos de evolución del pensamiento. Y puede ocurrir también que entidades que se han creado solo en la abstracción respondan luego a otras realidades. 

Al momento de establecer el objeto de estudio, previamente, hubo una consulta a los docentes en varias actividades y en esa búsqueda de cómo establecer la definición de qué es la Matemática, se determinó que ha sido objeto de transformaciones como producto de la evolución humana, del trabajo colectivo e individual de miles de personas conocidas como matemáticos, quienes visualizan a esta disciplina desde varias perspectivas: 

\begin{itemize}
    
     \color{blue} \item \color{black} Como ciencia, que busca dar respuesta a problemas de otras áreas del conocimiento mediante el razonamiento y la abstracción, que se ocupa de establecer relaciones formales entre objetos abstractos con el fin de explicar y comprender fenómenos de la naturaleza
    
    \color{blue} \item \color{black} Como fundamento, se ha establecido como base teórica de otras disciplinas por su abstracción, precisión, rigor lógico y su alcance para ofrecer una variedad de aplicaciones.
    
    \color{blue} \item \color{black} Como lenguaje, cuyas teorías establecidas se han comunicado de manera universal para las demás ciencias, con el fin de describir el Universo o por ser de interés para sí misma, donde se persigue entender el mundo, y el entorno matemático. 
    
    \color{blue} \item \color{black} Como una forma de expresión creativa: una forma de vida que activa esa pasión por entender los fenómenos de la naturaleza y modelarlos, que brinda un conjunto de técnicas para resolver problemas y que sirve como una herramienta organizativa, donde se describen, de forma rigurosa, aquellos comportamientos, fenómenos, estructuras e ideas que involucran cuantificación, clasificación, predicción, entre otros. 
    \end{itemize}

Con este proceso de reflexión suscitado entre los docentes de la Escuela de Matemática de la Universidad de Costa Rica, se comprende la Matemática desde sus áreas de estudio tales como: Álgebra, Análisis, Geometría, Modelización Matemática, Análisis Numérico, Lógica Matemática, Teoría de Números, Topología, Probabilidad y Ecuaciones Diferenciales.   Esta descripción de áreas no procura interpretarse como una forma de abordar la disciplina en forma separada, por el contrario, recientemente, se analizan situaciones de forma integrada considerando el aporte de las diferentes áreas del saber dentro de la Matemática; y, para este marco referencial, se denomina a cada área de las mencionadas anteriormente como identidad matemática.

Cada identidad matemática es parte de un conocimiento matemático que está representado en la Universidad de Costa Rica por profesionales que los caracteriza no solo por lo que saben, sino por lo que hacen. Al respecto, Mary Midgley expresa que “el conocimiento no es algo que tenemos y podemos comprar; es algo que hacemos y somos” [15]. Es decir, al hacer mención del conocimiento matemático que se produce en la Escuela de Matemática, también se hace referencia al progreso del colectivo que lo integra y con esto la oportunidad de llevar a cabo transformaciones con el paso del tiempo.



\subsubsection{Lo teórico-metodológico}

Se ha mencionado que la Matemática es una ciencia formal. Esta categoría que se le impregna provoca confusión en muchas ocasiones, de ahí que se requiera precisar algunos aspectos. Bunge en [9] comentó que: 

…la lógica y la matemática, por ocuparse de inventar entes formales y de establecer relaciones entre ellos, se llaman a menudo ciencias formales, precisamente porque sus objetos no son cosas ni procesos, sino, para emplear el lenguaje pictórico, formas en las que se puede verter un surtido ilimitado de contenidos, tanto fácticos como empíricos. Esto es, podemos establecer correspondencias entre esas formas (u objetos formales), por una parte, y cosas y procesos pertenecientes a cualquier nivel de la realidad por la otra.

El conocimiento que se establece dentro de la Matemática no es personal, es colectivo. Se logra con el intercambio de ideas y el trabajo colaborativo, al considerar la participación en diferentes ambientes de trabajo donde se plantean las investigaciones, los problemas (tanto los que se plantean dentro o fuera de la disciplina) a través de las reuniones científicas, las cuales se llevan a cabo por medio de: congresos, simposios, conferencias, coloquios, entre otras; en la modalidad de pequeñas reuniones donde se abordan temáticas de matemática especializada o en grandes eventos a nivel nacional e internacional.  

Este conocimiento matemático se produce mediante el pensamiento y la reflexión de las conjeturas que se establecen, el seguimiento que se lleva a cabo en la identificación de casos particulares, del estudio de ejemplos (particularmente cuando se procura resolver problemas que ya existen) y la importancia de seguir esa intuición, donde se genere un nuevo conocimiento, una nueva área de trabajo. Es ir brindando respuestas a algo que viene en proceso, o a una actividad práctica.En la Matemática siempre hay algo a lo cual responder o a lo que se está tratando de dar respuesta.   

La relación entre la teoría y la práctica es muy estrecha. Generalmente, la práctica es la propia investigación que se nutre de los congresos y de la interacción con otros especialistas, donde se logra producir nuevas ideas para un futuro trabajo de colaboración. La investigación de los demás es su práctica y nutre el trabajo a su vez nuevamente. Aunque, si se considera un entorno de prácticas profesionales en un sector productivo, se aplican modelos ya existentes, y, estos, a su vez, se pueden transformar en un espacio de generación de nuevos problemas que después pueden ser objeto de análisis por parte de un matemático. 

    Desde la filosofía de la matemática, Bunge en [9] aclara que:   
    El carácter matemático del conocimiento científico -esto es, el hecho de que es fundado, ordenado y coherente- es lo que lo hace racional. La racionalidad permite que el progreso científico se efectúe no sólo por la acumulación gradual de resultados, sino también por revoluciones. Las revoluciones científicas no son descubrimientos de nuevos hechos aislados, ni son perfeccionamientos en la exactitud de las observaciones, sino que consisten en la sustitución de hipótesis de gran alcance (principios) por nuevos axiomas, y en el reemplazo de teorías enteras por otros sistemas teóricos.  Sin embargo, semejantes revoluciones son a menudo provocadas por el descubrimiento de nuevos hechos de los que no dan cuenta las teorías anteriores, aunque a veces se encuentran en el proceso de comprobación de dichas teorías; y las nuevas teorías se tornan verificables en muchos casos, merced a la invención de nuevas técnicas de medición, de mayor precisión. 

\subsubsection{La finalidad}

Históricamente, la matemática se ha desarrollado como respuesta a las necesidades en el quehacer humano. De manera que el conocimiento matemático se produce para toda la humanidad, desde lo que se produce dentro de un ámbito académico, hasta en las experiencias de la cotidianidad, tal y como fueron comentados en el anterior marco referencial, principalmente en los
inicios de la historia.

Sin embargo, es importante resaltar las áreas o disciplinas que más matemática generan en la actualidad. Entre ellas se pueden citar: 

\begin{itemize}
    
    \color{blue} \item \color{black} El entorno socio económico. 

    
     \color{blue} \item \color{black} El entorno productivo (optimización, clasificación, big data, análisis numérico, etc.).  
     
    \color{blue} \item \color{black}  Las otras ciencias básicas: Física, Química, Biología, etc
    
     \color{blue} \item \color{black} El estudio de fenómenos sociodemográficos (modelación, etc.).
     
     \end{itemize}
     
La matemática se desarrolla, principalmente, por los matemáticos dedicados a la investigación, así como por profesionales de disciplinas afines. En las últimas décadas, la producción de conocimiento matemático, en instituciones financieras y del sector industrial, ha crecido significativamente desde el punto de vista de un quehacer interdisciplinario. Esto ha provocado que también se sumen otros profesionales de áreas afines, de manera más determinante, a la producción del conocimiento matemático.

A nivel nacional y a partir de la década de los años noventa, la producción matemática nacional ha estado básicamente desarrollada por dos centros de investigación (CIMPA y CIMM) de la Universidad de Costa Rica. Antes de la creación de estos centros de investigación, ya existía un desarrollo de trabajo en la matemática, pero gestionada directamente desde la Escuela de Matemática de la UCR. 

En un ámbito más general, el desarrollo de la matemática está, íntimamente, enlazado a universidades y centros de investigación, la mayoría de los cuales están ligados a universidades. Existen centros de investigación en algunos países que no están ligados a universidades (CNRS en Francia, por ejemplo), a pesar de que sus miembros sí deben ligarse a institutos de investigación. En los países del este de Europa, existen todavía las academias que son institutos de investigación no vinculados a Universidades. También existen algunos centros de investigación en salud o en distintas tecnologías que mantienen poca relación con instituciones de enseñanza superior, pero que contribuyen de manera significativa a la construcción de conocimiento matemático. 

No obstante, como se mencionó antes, la producción matemática de instituciones y empresas no ligadas a la academia representa una práctica emergente que ha crecido de manera significativa en las últimas décadas. 



\section{Definición de las áreas de conocimiento}


Se presentan las siguientes identidades matemáticas y no se pretende describirlas de manera exhaustiva ni tampoco expresar que son las únicas, pero son las que poseen relevancia o un mayor crecimiento en la Escuela de Matemática de la Universidad de Costa Rica. 

\begin{itemize}
    
     \color{blue} \item \color{black} \textbf{Álgebra:} El álgebra básica está fundamentada en la manipulación algebraica y resolución de ecuaciones y cálculos de naturaleza simples y sencillos. La evolución histórica del álgebra la transformó, más bien, en el estudio de las diferentes estructuras que llamamos algebraicas: espacios vectoriales, grupos, anillos, cuerpos y módulos en un primer abordaje. Asimismo, la interacción con otras áreas de la matemática, ha conducido a nuevas teorías y estructuras: teoría de representaciones, álgebra conmutativa, álgebra homológica, teoría de retículos, estructuras ordenadas, álgebras de Boole, de Heyting, de Lojasiewicz, teoría de formas cuadráticas, teoría de valuaciones, álgebra universal, algebra multilineal y teoría de álgebras sobre cuerpos. Cada una de ellas tiene interacciones con otras áreas de la matemática, por ejemplo: teoría de números, análisis, geometría y lógica. 
     
    \color{blue} \item \color{black}  \textbf{Análisis:} estudia operaciones que involucran conjuntos, funciones y espacios abstractos, en un ámbito de acción que, si bien, muchas veces es compartido con el álgebra, la naturaleza de esas operaciones es distinta de las manipulaciones algebraicas. Los objetos estudiados por el análisis y las interacciones entre estos, llevan implícito un carácter transfinito. Tal es el caso de la convergencia, la continuidad y otras propiedades topológicas. Son múltiples las subdivisiones que se pueden hacer del análisis, tanto, que es difícil en algunos casos establecer fronteras entre ellas o determinar cuál está contenida en cuál otra. Entre las subáreas del análisis se pueden mencionar el análisis real, dentro del cual podrían considerarse grandes áreas, como las ecuaciones diferenciales, la integración y teoría de la medida, la probabilidad y el análisis numérico. También podemos identificar el análisis complejo, análisis no estándar, análisis funcional y análisis armónico.
    
    \color{blue} \item \color{black} \textbf{Geometría:} es una de las principales ramas de la Matemática y combina la percepción visual con la lógica deductiva. En primer lugar, estudia las relaciones entre rectas y planos y entre curvas y superficies en el espacio, mediante métodos algebraicos con coordenadas y cálculo diferencial. Seguidamente, se abre a estructuras afines y proyectivas en dimensiones superiores, incorporando la teoría de variedades diferenciales y algebraicas. Comprende también el estudio de simetrías mediante grupos de Lie y sus acciones, amén de diversas nociones métricas: conexiones, curvatura e invariantes topológicos.
    
    \color{blue} \item \color{black} \textbf{Modelización Matemática:} describe algún fenómeno con una herramienta matemática. Existen diferentes tipos de modelaje, dentro de los que se destacan: dispersión, mediante ecuaciones diferenciales parciales y modelos espaciales; dinámica de población con ecuaciones diferenciales no lineales, modelos aleatorios con el uso de modelos estocásticos. Cabe mencionar que esta área se encuentra en constante desarrollo y se van implementando mejores formas de modelar. Es un area más interdisciplinaria (vista como herramienta), puesto que ayuda a contestar preguntas en campos fuera de la Matemática.
    
     \color{blue} \item \color{black} \textbf{Análisis Numérico:} área de la matemática que desarrolla métodos numéricos teóricos para resolver problemas prácticos. Involucra temas como el error numérico, eficiencia, complejidad, convergencia y estimaciones. Sirve como una herramienta aplicable a diferentes áreas.
     
    \color{blue} \item \color{black} \textbf{Lógica Matemática:} es el estudio de las interacciones entre el lenguaje formal y los objetos matemáticos.  Se divide en muchas áreas, algunas de las más importantes son: teoría de conjuntos, teoría de modelos, teoría de la demostración y recursividad algorítmica. La lógica matemática surge al inicio del siglo XX, a partir de los estudios del fundamento de las matemáticas; pero, posteriormente, llega a convertirse tanto en una herramienta, como en un área propia de las matemáticas. Tiene importantes interacciones con áreas como el álgebra, la teoría de números, la geometría, el análisis funcional y la informática teórica.
    
     \color{blue} \item \color{black} \textbf{Teoría de Números:} a un nivel elemental, estudia las propiedades de los números enteros primordialmente, entre las cuales se pueden mencionar las propiedades de los números primos y su distribución; y, la solubilidad de ecuaciones en los números enteros. Además, se estudian generalizaciones de este tipo de problemas, análogos de tipo algebraico y la representatividad de números en distintas formas. Es una de las áreas fundamentales de la matemática moderna, con aplicaciones muy importantes para la seguridad en las comunicaciones y transacciones a través de internet. Goza, actualmente, de mucho renombre, puesto que se ve representada en los mejores Departamentos de Matemática a nivel mundial, así como en las entregas de los premios más importantes en las matemáticas, como la medalla Fields y el premio Abel. En palabras de Gauss: “La reina de las Ciencias es la Matemática y la Teoría de los Números es la Reina de las Matemáticas”. 
     
    \color{blue} \item \color{black} \textbf{Topología:} es el estudio de las propiedades que posee un espacio que se preservan bajo deformaciones continuas. Uno de los objetivos centrales es la clasificación o diferenciación de los espacios topológicos, la cual se logra mediante el estudio de aquellos objetos que permanecen invariantes al transformar de manera continua un espacio topológico en otro. Ejemplos relevantes de espacios topológicos son los espacios métricos. Sin embargo, en la gran mayoría de casos, los aspectos propiamente geométricos se dejan para otras áreas de estudio. Esta área se concibe en el centro de la Matemática, es una herramienta que se usa en casi todas las áreas de esta ciencia y que ha tenido mucho éxito al incorporar en su haber otras áreas, como el álgebra y los aspectos diferenciales de la geometría.
    
     \color{blue} \item \color{black} \textbf{Probabilidad:} se expone su noción a partir de una lista de conceptos fundamentales tales como: medida de probabilidad, variable aleatoria, distribución y densidad de distribución, leyes de los grandes números, teorema del límite central, espacios de probabilidad, sigma-álgebra, procesos estocásticos, proceso de Markov, concepto de Martingala, generador, proceso de difusión, proceso de saltos, conceptos de cadena, caminata aleatoria, distintos tipos de convergencia. Dentro de la matemática, la probabilidad juega un rol integrador entre otras ramas porque aplica conceptos de muchas áreas, tales como: Análisis, Combinatoria, Álgebra Abstracta, Sistemas Dinámicos, entre otros. El lenguaje de la probabilidad y sus técnicas son también útiles en el enunciado y prueba de diversos resultados en esas áreas. Adicionalmente, muchas áreas de la matemática encuentran sus aplicaciones en otras ciencias a través de conceptos y resultados probabilísticos; ya que, se usa para representar y medir la incertidumbre de eventos reales. Los conceptos de variable aleatoria y proceso estocástico, propios de probabilidad, son la base teórica de las observaciones que usa la estadística.
     
    \color{blue} \item \color{black} \textbf{Ecuaciones Diferenciales:} son la expresión matemática de las leyes de la naturaleza. A pesar de tener un origen común, las ecuaciones diferenciales ordinarias y las ecuaciones en derivadas parciales, cuentan con métodos muy disímiles. Las ecuaciones ordinarias relacionan una o varias funciones desconocidas de una variable, así como sus derivadas. Las ecuaciones en derivadas parciales involucran funciones multivariables. En un primer momento las ecuaciones en derivadas parciales se centraron en el estudio de las ecuaciones lineales elípticas, parabólicas e hiperbólicas, para luego centrarse en las ecuaciones no lineales. Podría pensarse que la meta de la teoría de las ecuaciones diferenciales es, dada una ecuación, encontrar sus soluciones. Sin embargo, la mayoría de de las ecuaciones diferenciales no se puede resolver explícitamente. Por tanto, hacer un análisis cualitativo de las ecuaciones puede ser de mayor utilidad para el estudio de las soluciones.
    \end{itemize}

Para comprender la complejidad de lo que implica abstraer qué es la Matemática, el filósofo argentino Mario Bunge, al hacer alusión a la Lógica y Matemática, expresó que “a los lógicos y matemáticos no se les da objetos de estudio: ellos construyen sus propios objetos... La materia prima que emplean los lógicos y matemáticos no es fáctica sino ideal” ([9]). Se refuerza la noción de que es
una forma libre de la expresión humana, es parte de un gran espacio de inspiración; pero tiene sus reglas; es decir, un sistema establecido. Para los fines de este proceso, y, al tomar en consideración los aportes brindados por los matemáticos de la Escuela de Matemática, se estableció un objeto de estudio como cierre de la experiencia en la elaboración de este documento, lo cual implicó que no se limitara a solo asentir con lo expresado por otros matemáticos a nivel internacional. 

\subsubsection{Relación con otras disciplinas}

A partir de su tendencia a transformarse con el tiempo y de su espíritu de integración con otras áreas del conocimiento, la Matemática se ha cimentado al tomar como base no solo los avances que surgen dentro de ella, sino también los aportes de grandes áreas como: biología, demografía, microbiología, física, computación, filosofía, economía. 

Por otra parte, en función de ser una herramienta organizativa y que permita comprender mejor los fenómenos que se presentan, diferentes disciplinas emplean herramientas y teorías matemáticas para el avance de sus investigaciones, como: biología, ecología, dinámica social, ingeniería, física, química, computación, finanzas y arquitectura. 

Cabe destacar que la matemática no es una disciplina que ha dejado de poseer impacto. Por el contrario, recientemente se han fortalecido las herramientas existentes y se han introducido otras. Al especificar los campos de acción en los cuales incide la matemática en la investigación actual, se decidió recurrir a las identidades que fueron detalladas en el apartado del objeto de estudio, y, en las cuales, los docentes, según el área del conocimiento matemático expresaron que:

\begin{itemize}
    
    \color{blue} \item \color{black}En álgebra, hay tendencias teóricas y aplicadas. Con respecto a las aplicadas hay aportes en biología, física, química e ingeniería.En cuanto a las teóricas, con el advenimiento de las computadoras, esta tendencia cambió y se resalta la importancia de realizar experimentos computacionales para hacer conjeturas. Como tendencias específicas, se mencionan: álgebra homológica, álgebra no conmutativa, álgebra de Hopf, teoría de representaciones de grupos, problema inverso de Galois, colaboraciones en computación. 

    
     \color{blue} \item \color{black} En el análisis, se trabaja en la resolución numérica de ecuaciones diferenciales parciales, problemas inversos, ecuaciones no lineales, teoría de restricción y particiones polinomiales, análisis armónico. 
     
    \color{blue} \item \color{black}  En la modelización se destaca la oportunidad de ayudar al Ministerio de Salud en políticas de salud pública. Además, se ofrecen análisis para explicar cómo se propagan las enfermedades, o inclusive, se brindan insumos predictivos para elecciones presidenciales. 
    
     \color{blue} \item \color{black} En análisis numérico, por cada problema que surge en la industria, uno debe ser capaz de modelar numéricamente lo que esté pasando. En física, geología y meteorología, gran cantidad de modelos son ecuaciones que se resuelven numéricamente. 
     
     \color{blue} \item \color{black} La lógica matemática permite resolver problemas en otras disciplinas con nuevas herramientas. La teoría de modelos ha permitido avances en el álgebra y en la geometría algebraica que no se habían podido llevar a cabo por los propios especialistas. La teoría descriptiva de conjuntos puede verse como una solución a los problemas que yacen en la intersección de la lógica, la topología y el análisis. En informática teórica se aplica a la criptografía y a la teoría de complejidad algorítmica.
     
     \color{blue} \item \color{black} En teoría de números, hay vertientes de investigación que se concentran en el estudio de puentes entre distintas áreas dentro y fuera de la rama. Por mencionar: los avances recientes en el estudio de números primos, llegar a la demostración de la conjetura de los primos gemelos, demostración de la conjetura débil de Goldbach, y queda pendiente la conjetura fuerte.  
     
     \color{blue} \item \color{black} En topología, hay ramas de investigación que han consolidado aplicaciones en biología, ciencias de la computación, robótica, física; y hasta en juegos (rompecabezas) y arte. A lo interno de la matemática la topología se usa en casi todas las ramas, por ejemplo, la teoría de representaciones, la teoría de Lie, geometría diferencial y lógica. 
     
     \color{blue} \item \color{black} En probabilidad, hay necesidades puramente fundacionales (para extender los conceptos fundamentales) u obedece a aplicaciones, inspiradas por un problema que viene del área de las ciencias. 
     
     \color{blue} \item \color{black}  En ecuaciones diferenciales, el apartado es extenso y se citan algunas a continuación:
     
     \begin{itemize}
         
     
     \color{blue} \item \color{black} Investigación en imagen médica. Entender cómo se comporta la luz a través de los tejidos, trabajo que se lleva a cabo junto con las ecuaciones integrales. 
     
     \color{blue} \item \color{black} Métodos híbridos de imagen médica. En particular, las resonancias magnéticas, ultrasonidos, histografías, tomografía acústica computarizada (TAC).
     
     \color{blue} \item \color{black} En la epidemiología: análisis de vacunas, virus, control de plagas, análisis estructurales
     
     \color{blue} \item \color{black} En la Teoría de la Relatividad, Electromagnetismo, Mecánica Cuántica. 
     
     
      \color{blue} \item \color{black} En la Economía: modelos de flujo, capital.
      \end{itemize}
     
     \color{blue} \item \color{black} En geometría, hay aplicaciones en física, artes y arquitectura. También en la Teoría de Números, la Matemática Discreta, Química, Ingeniería, Modelo de poblaciones, entre otras.
     
    Deseamos enfatizar lo expresado por Bunge en [9]: podemos establecer correspondencias entre esas formas (u objetos formales), por una parte, y cosas y procesos pertenecientes a cualquier nivel de la realidad por la otra. Así es como la física, la química, la fisiología, la psicología, la economía, y las demás ciencias recurren a la matemática, empleándola como herramienta para realizar la más precisa reconstrucción de las complejas relaciones que se encuentran entre los hechos y sus diversos aspectos; dichas ciencias no identifican las formas ideales con los objetos concretos, sino que interpretan las primeras en términos de hechos y de experiencias (o, lo que es equivalente, formalizan enunciados fácticos).
     
     \end{itemize}